\documentclass[12pt, a4paper]{ctexart}
\usepackage[margin=2cm]{geometry}
\usepackage{libertine}

\usepackage{titlesec}
\usepackage{zhnumber}
\titleformat*{\section}{\Large\bfseries\raggedright}
\renewcommand{\thesection}{\zhnum{section}}
\renewcommand{\thesubsection}{\arabic{subsection}}

\setcounter{secnumdepth}{4}
\renewcommand{\paragraph}[1]{
    \par{\bfseries#1}\quad
    \vspace{2pt}
    \noindent
}
\newcommand{\question}[3]{
    \paragraph{题目描述} #1
    \paragraph{程序实现} \inputminted{java}{#2}
    \paragraph{运行结果} \inputminted{text}{#3}
}

\usepackage{enumitem}
\setlist[enumerate]{itemsep=2pt, parsep=0pt, partopsep=0pt, topsep=2pt}
\setlist[itemize]{itemsep=2pt, parsep=0pt, partopsep=0pt, topsep=2pt}
\linespread{1.2}

\usepackage{minted}
\setminted{
    breaklines=true, fontsize=\small, frame=lines,
    numbers=none, style=bw, tabsize=4
}
\usepackage{xparse}
\NewDocumentCommand{\mil}{ O{text} m }{ %
  \mintinline{#1}{#2} %
}

\usepackage{graphicx}
\usepackage{siunitx}
\usepackage{tikz}

\begin{document}

\pagestyle{plain}
\thispagestyle{empty}

\noindent
\begin{tabular*}{\textwidth}{l @{\extracolsep{\fill}} r @{\extracolsep{6pt}} l}
    \LARGE{\textbf{实验报告}} & Java 与 Web 开发 & \textit{Java and Web Development} \\
\end{tabular*}\\\\
\begin{tabular*}{\textwidth}{l l}
    \textbf{报告标题: } & \textbf{实验一} \\
    \textbf{学号: } & 19240212 \\
    \textbf{姓名: } & 华博文 \\
    \textbf{日期: } & 2026 年 3 月 20 日 \\
\end{tabular*}\\
\rule[2ex]{\textwidth}{2pt}

\section{实验环境}

\subsection{操作系统}

Ubuntu 24.04.3 LTS x86\_64，Linux 6.17.0-14-generic，AMD Ryzen 7 7700 (16) @ 5.3GHz，DDR5。

\subsection{编程工具}

NeoVim v0.12.0-dev，openjdk 21.0.10。

\subsection{其他工具}

\LaTeX{}。

\section{实验内容}

\subsection{第一部分：基本程序设计（必做题）}

\subsubsection{第 1 题}
\question{
    设计一个 Java 程序，完成：在屏幕上输出 ``九九乘法表''。
}{../src/A01/Solution.java}{../res/A01.out}

\subsubsection{第 2 题}
\question{
    设计一个 Java 程序，完成：从键盘上输入 \mil{n} $\left(1\leq\text{n}\leq10\right)$，\mil{t} $\left(4\leq\text{t}\leq8\right)$ 及输入一个字符 \mil{ch}，然后输出一个 ``树'' 形图形。
}{../src/A02/Solution.java}{../res/A02.out}

\subsubsection{第 3 题}
\question{
    设计一个 Java 程序，完成：已知四个 \mil[java]{byte} 类型的变量 \mil{b1}、\mil{b2}、\mil{b3}、\mil{b4}，它们的值分别是（例如）\mil{0xaa}、\mil{0xbb}、\mil{0xcc}、\mil{0xdd}；现要求将它们拼接成一个 \mil[java]{int} 类型的值（$4$ 个字节），按次序 $b4b3b2b1$ 拼装后，放到 \mil[java]{int i} 变量中，最后打出 \mil[java]{i} 的值。
}{../src/A03/Solution.java}{../res/A03.out}

\subsubsection{第 4 题}
\question{
    设计一个 Java 程序，完成：$1+3+5+\dots+n$ 之和。若 $n$ 是偶数，只加到 $n-1$。若 $n$ 是奇数，则正好加到 $n$ 为止。
}{../src/A04/Solution.java}{../res/A04.out}

\subsubsection{第 5 题}
\question{
    设计一个 Java 程序，完成：$i+\left(i+1\right)+\left(i+2\right)+\dots+j$ 之和。
}{../src/A05/Solution.java}{../res/A05.out}

\subsubsection{第 6 题}
\question{
    设计一个 Java 程序，完成：$i+\left(i+k\right)+\left(i+2k\right)+\left(i+3k\right)+\dots+j$ 之和。
}{../src/A06/Solution.java}{../res/A06.out}

\subsubsection{第 7 题}
\question{
    设计一个 Java 程序，完成：输出 $i\sim j$ 之间所有的质数，每行 $k$ 个。
}{../src/A07/Solution.java}{../res/A07.out}

\subsubsection{第 8 题}
\question{
    设计一个 Java 程序，完成：输出前 $n$ 个质数，每行 $k$ 个。
}{../src/A08/Solution.java}{../res/A08.out}

\subsubsection{第 9 题}
\question{
    设计一个 Java 程序，完成：输入正整数 $n$，计算 $$1-\frac{1}{2}+\frac{1}{3}-\frac{1}{4}+\frac{1}{5}-\frac{1}{6}+\dots\pm\frac{1}{n}$$
}{../src/A09/Solution.java}{../res/A09.out}

\subsection{第二部分：竞赛预备基础题（提高题）}

\subsubsection{第 1 题}
\question{
    设计一个 Java 程序，完成：输入两个正整数 $a$ 和 $b$，输出它们的最大公约数。
}{../src/B01/Solution.java}{../res/B01.out}

\subsubsection{第 2 题}
\question{
    设计一个 Java 程序，完成：输出所有的形如 $\overline{aabb}$ 的四位平方数 （即前两位数字相等、后两位数字相等、且是平方数）。
}{../src/B02/Solution.java}{../res/B02.out}

\subsubsection{第 3 题}
\question{
    设计一个 Java 程序，完成：大数计算，求阶乘之和。输入整数 $n$ $\left(n<10^6\right)$，求出 $s=1!+2!+\dots+n!$ 这个很大的和的末 $6$ 位。
}{../src/B03/Solution.java}{../res/B03.out}

\subsubsection{第 4 题}
\question{
    设计一个 Java 程序，完成：队伍人数。一支队伍，若 $3$ 个一排，则队尾有 $a$ 人 $\left(a<3\right)$，若 $5$ 个一排，则队尾有 $b$ 个 $\left(b<5\right)$，若 $7$ 个一排，则队尾有 $c$ 人 $\left(c<7\right)$。现从键盘上输入 $a$、$b$、$c$ 的值，输出满足该要求的总人数的最小值，或者报告 ``无解''。已知总人数大于等于 $10$ 但小于 $100$。
}{../src/B04/Solution.java}{../res/B04.out}

\subsubsection{第 5 题}
\question{
    设计一个 Java 程序，完成：用 $1,2,3,\dots,9$ 组成 $3$ 个 $3$ 位数：$\overline{abc}$、$\overline{def}$、$\overline{ghi}$。要求每个数字恰好使用了一次，且 $\overline{abc}:\overline{def}:\overline{ghi}=1:2:3$。输出所有的解。
}{../src/B05/Solution.java}{../res/B05.out}

\subsubsection{第 6 题}
\question{
    设计一个 Java 程序，完成：输入整数 $n$，输出对应的 ``康特数''。康特数的排列如下：
    \begin{figure}[h]
        \centering\begin{tikzpicture}[
    node distance=2cm,
    font=\small,
    every node/.style={
        minimum size=0.8cm,
        align=center
    },
    arrowstyle/.style={red, ->, thick}
]

\node (n11) at (0, 4.8) {$1/1$};
\node (n12) at (1.6, 4.8) {$1/2$};
\node (n13) at (3.2, 4.8) {$1/3$};
\node (n14) at (4.8, 4.8) {$1/4$};
\node (n15) at (6.4, 4.8) {$1/5$};

\node (n21) at (0, 3.6) {$2/1$};
\node (n22) at (1.6, 3.6) {$2/2$};
\node (n23) at (3.2, 3.6) {$2/3$};
\node (n24) at (4.8, 3.6) {$2/4$};

\node (n31) at (0, 2.4) {$3/1$};
\node (n32) at (1.6, 2.4) {$3/2$};
\node (n33) at (3.2, 2.4) {$3/3$};

\node (n41) at (0, 1.2) {$4/1$};
\node (n42) at (1.6, 1.2) {$4/2$};

\node (n51) at (0, 0) {$5/1$};

\draw[arrowstyle] (n11) -- (n12);
\draw[arrowstyle] (n12) -- (n21);

\draw[arrowstyle] (n21) -- (n31);
\draw[arrowstyle] (n31) -- (n22);
\draw[arrowstyle] (n22) -- (n13);

\draw[arrowstyle] (n13) -- (n14);
\draw[arrowstyle] (n14) -- (n23);
\draw[arrowstyle] (n23) -- (n32);
\draw[arrowstyle] (n32) -- (n41);

\draw[arrowstyle] (n41) -- (n51);
\draw[arrowstyle] (n51) -- (n42);
\draw[arrowstyle] (n42) -- (n33);
\draw[arrowstyle] (n33) -- (n24);
\draw[arrowstyle] (n24) -- (n15);

\node (n16) at (8, 4.8) {$\cdots$};
\draw[arrowstyle] (n15) -- (n16);

\end{tikzpicture}
    \end{figure}
}{../src/B06/Solution.java}{../res/B06.out}

\subsubsection{第 7 题}
\question{
    设计一个 Java 程序，完成：有一群海盗（不多于 $20$ 人）在船上比拼酒量。过程如下：打开一瓶酒，所有在场的人平分喝下，有几个人倒下了；再打开一瓶酒平分，又有倒下的；再次重复，直到开了第 $4$ 瓶酒，坐着的已经所剩无几，海盗船长也在其中。当第 $4$ 瓶酒平分喝下后，大家都倒下了。
    
    等船长醒来，发现海盗船搁浅了。他在航海日志中写道：``昨天，我正好喝了一瓶...奉劝大家，开船不喝酒，喝酒不开船。'' 请你根据这些信息，推断开始有多少人，每一轮喝下来还剩多少人。如果有多个可能的答案，请列出所有答案，每个答案占一行。
}{../src/B07/Solution.java}{../res/B07.out}

\end{document}
